\documentclass[12pt]{article}

\usepackage[margin=1in]{geometry}
\usepackage{amsmath}
\usepackage{booktabs}
\usepackage{setspace}
\usepackage{hyperref}

\hypersetup{
  colorlinks=true,
  linkcolor=blue,
  citecolor=blue,
  urlcolor=blue
}

\title{Regular Labor:\\
Chilean Prune Shocks, Digestive Health, and Earnings}

\author{
  Serdar Ozkan\thanks{All estimates in this draft are illustrative. Strategic prune reserve policy simulations are, at present, aspirational.}
  \and
  Codex\thanks{Correspondence welcome, especially from readers with strong priors on the exclusion restriction.}
}

\date{\today}

\begin{document}

\maketitle

\begin{abstract}
\noindent
This paper studies how short-run disruptions to digestive health affect labor-market outcomes. We use weather shocks in Chilean prune-producing regions as an instrument for Danish prune supply, exploiting Denmark's reliance on Chile as its largest direct source of dried prunes. A one-standard-deviation adverse weather shock in Chile reduces Danish prune imports by 4.6 percent and raises retail prune prices by 2.5 percent. Under the exclusion restriction that Chilean growing-season weather affects Danish earnings only through prune availability, the IV estimates imply that a 10 percent increase in prune prices raises constipation-related primary-care contacts by 3.2 percent and reduces monthly earnings by 0.12 percent. The earnings effects operate through small increases in sickness absence and reductions in paid hours, and are concentrated among older workers, workers with prior constipation-related health contacts, and occupations with limited schedule flexibility. Results are similar, though less precise, when we instrument for Danish prune availability using weather shocks in California's Central Valley, the largest global prune-producing region. The findings suggest that a strategic prune reserve would be a low-cost, if unusually specific, form of labor-market stabilization policy.
\end{abstract}

\noindent
\textbf{Keywords:} digestive health, labor supply, earnings, weather shocks, instrumental variables, prunes

\noindent
\textbf{JEL Codes:} I12, J22, J31, Q54

\onehalfspacing

\iffalse
\section{Introduction}

Mild health frictions are common, transitory, and difficult to study. Their effects on labor-market outcomes may nevertheless be measurable when workers face rigid schedules, limited control over breaks, or jobs that require physical presence. This paper studies one such friction: constipation.

We exploit exogenous variation in Danish prune availability generated by weather shocks in Chilean prune-producing regions. The design uses Chilean growing-season weather as an instrument for Danish prune prices and imports. The first stage arises because Denmark relies heavily on Chile as a direct source of dried prunes. The exclusion restriction requires that Chilean weather affect Danish earnings only through its effect on prune supply, rather than through other channels linking Chilean agriculture to Danish labor-market conditions.

The estimates are intentionally small. A 10 percent increase in prune prices raises constipation-related primary-care contacts by 3.2 percent and reduces monthly earnings by 0.12 percent. These effects are concentrated among older workers, workers with prior constipation-related health contacts, and workers in low-flexibility occupations. The results therefore speak less to large productivity losses than to a narrow but concrete margin of labor-market fragility.

\section{Empirical Strategy}

Let $i$ index workers, $m$ months, and $t$ years. The baseline first stage relates Danish prune prices to weather shocks in Chile:

\begin{equation}
  \log P_{mt}^{\text{prune}} =
  \pi Z_{mt}^{\text{Chile weather}} + \gamma_m + \delta_t + X_{mt}'\lambda + u_{mt},
\end{equation}

\noindent
where $P_{mt}^{\text{prune}}$ is the Danish retail prune price, $Z_{mt}^{\text{Chile weather}}$ summarizes adverse growing-season weather in Chilean prune-producing regions, $\gamma_m$ are calendar-month fixed effects, and $\delta_t$ are year fixed effects.

The second stage estimates the effect of prune prices on worker outcomes:

\begin{equation}
  Y_{imt} =
  \beta \widehat{\log P}_{mt}^{\text{prune}} + \gamma_m + \delta_t + X_{imt}'\theta + \varepsilon_{imt},
\end{equation}

\noindent
where $Y_{imt}$ denotes earnings, paid hours, sickness absence, or constipation-related health contacts. The coefficient $\beta$ captures the causal effect of prune prices for workers whose digestive-health inputs respond to imported prune supply.

\section{Illustrative Results}

\begin{table}[h!]
\centering
\caption{Illustrative Effects of Prune Price Shocks}
\begin{tabular}{lcc}
\toprule
Outcome & Effect of 10\% Price Increase & Unit \\
\midrule
Constipation-related primary-care contacts & 3.2 & Percent \\
Monthly earnings & -0.12 & Percent \\
Monthly earnings & -36 & DKK \\
Sickness absence & 0.004 & Days \\
Paid hours & -0.018 & Hours \\
\bottomrule
\end{tabular}
\end{table}

\section{Policy Implication}

The estimates imply that constipation is not a first-order determinant of aggregate labor supply. They do, however, identify a small and preventable source of worker discomfort that matters most when jobs offer little flexibility. A strategic prune reserve would therefore be a low-cost, if unusually specific, form of labor-market stabilization policy.

\fi
\end{document}
